\section{Modificación de datos}

Para no alterar observaciones reales de INUMET, las cinco modificaciones se realizaron sobre registros sintéticos claramente identificables. En cada caso se documentó el estado anterior, la operación y el estado posterior, tal como exige la consigna.

\subsection{INSERT de una nueva medición}

Se verificó primero que no existiera la observación identificada como \texttt{PRUEBA\_UTEC}.

\begin{lstlisting}[language=CQL]
SELECT *
FROM mediciones_por_estacion_mes
WHERE estacion_id = 'PRUEBA_UTEC'
  AND anio = 2025
  AND mes = 1
  AND fecha_hora = '2025-01-15 12:30:00';
\end{lstlisting}

\begin{figure}[H]
    \centering
    \includegraphics[width=\figwidth]{imagenes/capturas/08_insert_antes.png}
    \caption{Estado anterior al INSERT: no existía una fila con la clave primaria de \texttt{PRUEBA\_UTEC}.}
    \label{fig:insert-antes}
\end{figure}

La nueva observación se insertó con temperatura 24,5, humedad 65 y precipitación 1,2:

\begin{lstlisting}[language=CQL]
INSERT INTO mediciones_por_estacion_mes (
  estacion_id, anio, mes, fecha_hora,
  temp_aire, hum_relativa, precip_horario
)
VALUES (
  'PRUEBA_UTEC', 2025, 1,
  '2025-01-15 12:30:00',
  24.5, 65, 1.2
);
\end{lstlisting}

\begin{figure}[H]
    \centering
    \includegraphics[width=0.78\linewidth]{imagenes/capturas/09_insert_operacion.png}
    \caption{Ejecución del INSERT del registro sintético \texttt{PRUEBA\_UTEC}.}
    \label{fig:insert-op}
\end{figure}

\begin{figure}[H]
    \centering
    \includegraphics[width=\figwidth]{imagenes/capturas/10_insert_despues.png}
    \caption{Estado posterior al INSERT: la observación sintética aparece almacenada.}
    \label{fig:insert-despues}
\end{figure}

\subsection{UPDATE de la temperatura}

Se modificó únicamente \texttt{temp\_aire}, pasando de 24,5 a 27,0.

\begin{figure}[H]
    \centering
    \includegraphics[width=\figwidth]{imagenes/capturas/11_update_antes.png}
    \caption{Estado anterior al UPDATE: \texttt{temp\_aire = 24.5}.}
    \label{fig:update-antes}
\end{figure}

\begin{lstlisting}[language=CQL]
UPDATE mediciones_por_estacion_mes
SET temp_aire = 27.0
WHERE estacion_id = 'PRUEBA_UTEC'
  AND anio = 2025
  AND mes = 1
  AND fecha_hora = '2025-01-15 12:30:00';
\end{lstlisting}

\begin{figure}[H]
    \centering
    \includegraphics[width=\figwidth]{imagenes/capturas/12_update_despues.png}
    \caption{Estado posterior al UPDATE: la temperatura pasó a 27,0.}
    \label{fig:update-despues}
\end{figure}

\subsection{Sobrescritura mediante INSERT sobre la misma PRIMARY KEY}

A continuación se ejecutó un nuevo \texttt{INSERT} con la misma clave primaria del registro sintético, pero con nuevos valores. Cassandra no comprueba por defecto si la fila ya existía: una escritura con la misma clave actualiza los valores correspondientes, comportamiento de tipo \emph{upsert}. \parencite{cassandra_cql}

\begin{lstlisting}[language=CQL]
INSERT INTO mediciones_por_estacion_mes (
  estacion_id, anio, mes, fecha_hora,
  temp_aire, hum_relativa, precip_horario
)
VALUES (
  'PRUEBA_UTEC', 2025, 1,
  '2025-01-15 12:30:00',
  18.0, 80, 5.5
);
\end{lstlisting}

El estado anterior corresponde a la Figura~\ref{fig:update-despues}, con temperatura 27,0, humedad 65 y precipitación 1,2. La Figura~\ref{fig:upsert-despues} confirma que la misma clave contiene ahora los valores 18,0, 80 y 5,5.

\begin{figure}[H]
    \centering
    \includegraphics[width=\figwidth]{imagenes/capturas/13_upsert_despues.png}
    \caption{Estado posterior a la sobrescritura: la misma PRIMARY KEY conserva una única fila con los nuevos valores.}
    \label{fig:upsert-despues}
\end{figure}

\subsection{INSERT USING TTL}

Se utilizó el registro sintético \texttt{UTEC\_TTL} para demostrar expiración temporal. Antes de la operación no existía ninguna fila con esa clave.

\begin{figure}[H]
    \centering
    \includegraphics[width=\figwidth]{imagenes/capturas/14_ttl_antes.png}
    \caption{Estado anterior al INSERT con TTL: la fila \texttt{UTEC\_TTL} no existía.}
    \label{fig:ttl-antes}
\end{figure}

\begin{lstlisting}[language=CQL]
INSERT INTO mediciones_por_estacion_mes (
  estacion_id, anio, mes, fecha_hora,
  temp_aire, hum_relativa, precip_horario
)
VALUES (
  'UTEC_TTL', 2025, 1,
  '2025-01-15 13:00:00',
  22.0, 70, 0.5
)
USING TTL 300;
\end{lstlisting}

\texttt{USING TTL 300} asignó un tiempo de vida de 300 segundos a los valores escritos. La función \texttt{TTL(temp\_aire)} permitió consultar el tiempo restante; en la captura se observan 282 segundos.

\begin{figure}[H]
    \centering
    \includegraphics[width=\figwidth]{imagenes/capturas/15_ttl_despues.png}
    \caption{Estado posterior al INSERT con TTL. La función \texttt{TTL(temp\_aire)} devuelve 282 segundos restantes.}
    \label{fig:ttl-despues}
\end{figure}

\subsection{DELETE}

Finalmente se eliminó el registro sintético \texttt{PRUEBA\_UTEC}. Antes del borrado todavía se observaban los valores establecidos por el upsert.

\begin{figure}[H]
    \centering
    \includegraphics[width=\figwidth]{imagenes/capturas/16_delete_antes.png}
    \caption{Estado anterior al DELETE del registro \texttt{PRUEBA\_UTEC}.}
    \label{fig:delete-antes}
\end{figure}

\begin{lstlisting}[language=CQL]
DELETE FROM mediciones_por_estacion_mes
WHERE estacion_id = 'PRUEBA_UTEC'
  AND anio = 2025
  AND mes = 1
  AND fecha_hora = '2025-01-15 12:30:00';
\end{lstlisting}

\begin{figure}[H]
    \centering
    \includegraphics[width=\figwidth]{imagenes/capturas/17_delete_despues.png}
    \caption{Estado posterior al DELETE: la consulta ya no devuelve la fila eliminada.}
    \label{fig:delete-despues}
\end{figure}
